\section{Scoping Audit Protocol and Empirical Test Bench}
\label{sec:protocol}

We conducted an empirical scoping audit in accordance with the PRISMA Extension for Scoping Reviews (PRISMA-ScR) guidelines \cite{tricco2018} and established scoping methodology \cite{levac2010}. The audit protocol was specified prior to full-text extraction to govern retrieval strings, eligibility thresholds, multi-pass screening, and inter-rater deduplication.

\subsection{Information Sources and Retrieval Strategy}
We retrieved candidate literature across two complementary discovery channels covering January 1, 2023, through March 1, 2026, with rolling monitoring through July 2026 for emergent preprints. First, we conducted programmatic searches via the OpenAlex API, which indexes records across IEEE Xplore, ScienceDirect, SpringerLink, Wiley, and MDPI. Using four targeted title queries---$Q_1$ (\texttt{"physics-constrained" AND "wind"}), $Q_2$ (\texttt{"physics-informed" AND "wind power"}), $Q_3$ (\texttt{"physics-guided" AND "wind"}), and $Q_4$ (\texttt{"physics-embedded" AND "wind power"})---we identified 72 candidate records. Second, we conducted targeted searches on Google Scholar, arXiv, and preprint servers to capture early-access journal articles and working papers, yielding 17 additional records ($N=89$ total). Candidate records were cross-checked against Scopus and IEEE Xplore to verify coverage of major power systems and AI venues. All search strings, raw records, and exclusion criteria are documented in our open supplementary ledger. Furthermore, a targeted sensitivity search expanding retrieval to title, abstract, and author keywords across Scopus and IEEE Xplore (encompassing broader descriptors such as \texttt{"knowledge-guided"}, \texttt{"hybrid"}, and \texttt{"domain-informed"}) identified 14 additional papers. Full-text review confirmed that none documented an operational failure boundary or negative trade-off, specifically verifying that the 0.0\% negative-harm prevalence finding is robust to broader search scope. We note that this targeted 14-study query served strictly to stress-test negative trade-off reporting; exhaustive five-dimensional extraction ($\mathcal{D}_1$--$\mathcal{D}_5$) and inter-rater reliability coding were conducted on the core PRISMA-ScR scoping cohort of 46 studies.

\begin{figure}[!t]
\centering
\includegraphics[width=\columnwidth]{prisma_flow_diagram.pdf}
\caption{PRISMA-ScR Flow Diagram Illustrating the Scoping Corpus Derivation ($89 \rightarrow 72 \rightarrow 53 \rightarrow 48 \rightarrow 46$). Deduplication merged 5 preprint/journal pairs, while full-text eligibility screening excluded 2 studies, leaving 46 eligible full-text studies.}
\label{fig:prisma_flow}
\end{figure}

\subsection{Screening, Deduplication, and Eligibility Criteria}
Candidate records underwent structured multi-phase screening following the PRISMA-ScR workflow (Figure~\ref{fig:prisma_flow}). Cross-source deduplication first eliminated 17 duplicate records from the initial 89 items, leaving 72 unique entries for title and abstract evaluation. Screening at this stage excluded 19 records with documented rationale: nine addressed non-forecasting tasks such as turbine structural health monitoring, blade crack detection, or aeroelastic fatigue modeling; nine focused on real-time blade pitch or active power control rather than multi-step power forecasting; and one contained no physics-based modeling components. This filtering yielded 53 included candidate studies sought for full-text retrieval.

Because authors frequently release preliminary working papers on preprint servers (such as arXiv or SSRN) prior to formal journal publication, treating preprints and published articles independently would artificially inflate corpus statistics. We enforced a strict study-level deduplication rule: whenever a preprint and a peer-reviewed journal paper shared identical authorship, empirical datasets, and mathematical formulations, they were merged into a single study-level record retaining the peer-reviewed journal citation. Five paired duplicate sets were identified and consolidated (e.g., \textit{Applied Energy} 2025 paired with SSRN 2024, and \textit{Electric Power Systems Research} 2026 paired with SSRN 2026), reducing the candidate cohort from 53 to 48 unique study-level entries.

The remaining 48 full texts were independently evaluated against three explicit scoping eligibility criteria: (1) primary focus on short-term or ultra-short-term wind power forecasting; (2) mathematical formulation and computational execution of an explicit physical constraint, governing equation, or domain envelope; and (3) empirical validation on real-world or simulated wind power generation time series. Two studies failed to meet these criteria: study S26 was excluded because it forecasted atmospheric wind velocity exclusively without an explicit turbine power constraint, while study S42 was excluded because it implemented heuristic feature clustering without domain physical equations. This systematic filtering established our final representative scoping corpus of \textbf{$N=46$ unique scoping studies}, comprising 38 peer-reviewed journal articles and eight representative preprints (including UniWind \cite{s27_uniwind2026}).

\subsection{Multi-Rater Reliability Benchmarking}
To assess the empirical stability of the Schema v4 taxonomy before full ledger extraction, three independent raters (Coder A, Coder B, and Coder C) double-blindly coded a stratified representative subset of 20 candidate studies across all five dimensions. The pilot benchmarking achieved substantial to perfect inter-rater concordance across taxonomy facets. Study eligibility screening achieved unanimous consensus (Fleiss $\kappa = 1.000$), correctly isolating negative controls S26 and S42. Parameter provenance ($\mathcal{D}_1$) demonstrated perfect agreement ($\kappa = 1.000$ across both total and included corpora), while application stage ($\mathcal{D}_2$) yielded near-perfect agreement ($\kappa = 0.950$ overall, $\kappa = 0.938$ on included studies). Inference telemetry requirements ($\mathcal{D}_3$) achieved substantial agreement ($\kappa = 0.743$ overall, $\kappa = 0.646$ on included studies).

In contrast, implementation fidelity ($\mathcal{D}_4$) exhibited moderate concordance prior to consensus arbitration ($\kappa = 0.716$ overall in the stratified $N=20$ pilot with negative controls, but $\kappa = 0.341$ when evaluated strictly on included studies). This reduction reflects the classical prevalence paradox described by Feinstein and Cicchetti: because 92.6\% of candidate studies omitted public code, expected chance agreement $P_e$ was high ($86.1\%$), which depresses Fleiss' $\kappa$ despite a raw observed agreement of $P_o = 90.7\%$. Evaluating static manuscripts without code repositories leaves coders to interpret ambiguous narrative descriptions regarding whether a constraint represents an active computational component or conceptual motivation. This divergence indicates that software implementation fidelity cannot be reliably inferred from published prose alone, motivating direct computational graph inspection. For ablation semantics ($\mathcal{D}_5$), raters achieved strong concordance ($\kappa = 0.851$ overall, $\kappa = 0.818$ on included studies). On harm reporting, all three raters unanimously coded zero instances of harm documentation across the 18 included studies ($0/18$); because marginal variance was zero ($P_e = 1.0$), Cohen's $\kappa$ is mathematically indeterminate ($0/0$) and is formally reported as 100\% observed agreement ($P_o = 1.000$). Following pilot calibration and consensus arbitration, full-text extraction across the complete 46-study corpus was conducted under Schema v4 at the study level (coding each study's primary operational constraint mechanism alongside verbatim textual citations and equation coordinates), with primary extraction executed by J.L. and cross-verification of key quantitative parameters conducted by J.D.

\subsection{Software Audit Protocol (PCT \& PINN-UQ)}
The code availability audit followed a standardized five-step software search protocol examining full-text bodies, footnotes, data availability statements, publisher repositories, and author profiles. Across the 46-study corpus, only two publicly accessible software repositories were identified. One repository, PINN-UQ (S10, \textit{Engineering Applications of Artificial Intelligence} 2025), was audited and verified to faithfully implement its specified algebraic power equations ($P = g T \omega$ and $P = \frac{1}{2} C_p \rho A v^3$ via \texttt{sn.Tie}); its computational graph strictly executes the published algebraic electromechanical relationships without unstated differential substitutions, earning classification as \textit{verified-aligned}. The remaining repository, the Physics-Constrained Transformer (PCT, S02, \textit{Scientific Reports} 2026), provided open PyTorch scripts (\texttt{train\_clean.py} and \texttt{tst/loss.py}) which we audited line-by-line for parity between published loss formulations and computational graph execution, as detailed in Section~\ref{sec:code_audit}.

\subsection{Downstream Operational Dispatch and Settlement Test Bench}
To ground the empirical audit in real-world power system operations across diverse meteorological and topographic terrains, we established an operational dispatch test bench across three commercial UK wind farms with 10-minute SCADA telemetry. The \textit{Hill of Towie Wind Farm} (48.3~MW rated capacity, 21 Siemens SWT-2.3-82 turbines) features complex upland terrain and serves as an independent spatiotemporal extrapolation benchmark (7,599 test origins, 45,594 delivery steps across full-year 2017), with data released by RES on behalf of TRIG openly available on Zenodo under CC-BY-4.0 \cite{zenodo_hilloftowie}. The \textit{Kelmarsh Wind Farm} (12.3~MW rated capacity, six Senvion MM92 turbines) is situated in flat agricultural terrain and serves as our primary development benchmark (8,212 test origins, 49,272 delivery steps across full-year 2017) \cite{zenodo_kelmarsh}. The \textit{Penmanshiel Wind Farm} (28.7~MW rated capacity, 14 Senvion MM82 turbines, WT01--WT15 with WT03 absent) is located in the Scottish Borders \cite{zenodo_penmanshiel}, providing an independent upland setting with 8,440 test origins configured with a deliberately restricted training partition ($N_{\text{train}} = 1,583$) to evaluate model behavior under sample-restricted regimes. We note that because Penmanshiel exhibits distinct turbine models and terrain compared to Kelmarsh and Hill of Towie, cross-site comparisons inherently capture both site-specific topography and sample size differences; within-site controlled ablations remain essential to decouple these factors.

The empirical test bench operates under a prospective, frozen evaluation protocol. The forecasting pipeline adopts an input lookback of $L=36$ (6 hours) and a forecast horizon of $H=36$ (6 hours), with downstream economic dispatch evaluated over the next-hour delivery window $H_{\text{disp}} = \{1, \dots, 6\}$ and an origination stride of $S=6$ steps to prevent overlapping settlement intervals. The core model suite spans unconstrained DLinear, Bounds-Constrained DLinear ($\lambda=0.01$), Empirical Power Curve Constrained DLinear ($\lambda=0.01$), and Conditional Distribution Prior DLinear ($\lambda=0.01$), benchmarked alongside classical Ridge regression and Persistence baselines across five random initialization seeds. To verify whether empirical findings extend beyond linear representations, we additionally evaluate a state-of-the-art multi-patch channel-independent Transformer backbone (\textbf{PatchTST} \cite{nie2023patchtst}, patch length 6, stride 4, $d_{\text{model}}=64$, 2-layer encoder with RevIN) across all three wind farms, four constraint classes, and five random seeds (60 trained models) as an architectural sensitivity baseline under univariate power input and fixed regularization ($\lambda=0.01$).

Downstream decision calibration is implemented via a post-hoc quantile offset $\delta_h$ frozen strictly from the independent 2016 calibration year ($q = 1/(r+1) = 1/3$ for penalty ratio $r=2$) with zero access to 2017 test telemetry or market prices. Real-world physical settlement is evaluated against official 2017 UK Elexon BSC System Buy Prices (mean \pounds 44.27/MWh, peak \pounds 1,509.80/MWh, with 68 negative price periods down to -\pounds 73.15/MWh). Cashflows account for 30-minute physical netting:
\begin{equation}
E_{SP} = \sum_{t \in SP} (p_t - y_t) \cdot \frac{P_{\text{rated}}}{6000} \quad [\text{MWh}],
\label{eq:bsc_netting}
\end{equation}
generating single-sided shortfall cashouts and two-sided imbalance-settlement cashflows as a decision-risk diagnostic, validated via 2,000-replication 7-day circular block paired bootstrapping.
